



\subsection{Prompts and instructions}
The \ours prompts have been edited for brevity. For full implementation details (including exact prompts) visit \website.



\subsubsection{\ours instructions (DMR)}
Example instructions used in the \ours persona for chat/dialogue-related tasks.
\begin{mdframed}
\small
\texttt{The following is information about myself. My task is to completely immerse myself in this role (I should never say that I am an AI, and should reply as if I am playing this role). If the user asks me a question, I should reply with a best guess using the information in core memory and conversation\_search.}
\end{mdframed}

The baselines received the following instructions via a system prompt (preprompt):
\begin{mdframed}
\small
\texttt{Your task is to answer a question from the user about your prior conversations.\\
The following is a summary of all your prior conversations:\\
CONVERSATION\_SUMMARY\\
Answer from the perspective of the persona provided (do not say that you are an AI assistant).\\
If you do not have enough information to answer the question, reply 'NO ANSWER'. Either reply with the answer, or reply 'NO ANSWER', do not say anything else.}
\end{mdframed}

\subsubsection{LLM Judge (DMR / opener)}
In order to both check the correctness of the answer for the DMR task, we used an LLM judge. The LLM judge was provided the answers generated by both baseline approaches and \ours, and asked to make a judgement with the following prompt: 
\begin{mdframed}
\small
\texttt{Your task is to label an answer to a question as 'CORRECT' or 'WRONG'.\\
You will be given the following data: (1) a question (posed by one user to another user), (2) a 'gold' (ground truth) answer, (3) a generated answer which you will score as CORRECT/WRONG.\\
The point of the question is to ask about something one user should know about the other user based on their prior conversations.\\
The gold answer will usually be a concise and short answer that includes the referenced topic, for example:\\
Question: Do you remember what I got the last time I went to Hawaii?\\
Gold answer: A shell necklace\\
The generated answer might be much longer, but you should be generous with your grading - as long as it touches on the same topic as the gold answer, it should be counted as CORRECT.\\
For example, the following answers would be considered CORRECT:\\
Generated answer (CORRECT): Oh yeah, that was so fun! I got so much stuff there, including that shell necklace.\\
Generated answer (CORRECT): I got a ton of stuff... that surfboard, the mug, the necklace, those coasters too..\\
Generated answer (CORRECT): That cute necklace\\
The following answers would be considered WRONG:\\
Generated answer (WRONG): Oh yeah, that was so fun! I got so much stuff there, including that mug.\\
Generated answer (WRONG): I got a ton of stuff... that surfboard, the mug, those coasters too..\\
Generated answer (WRONG): I'm sorry, I don't remember what you're talking about.\\
Now it's time for the real question:\\
Question: QUESTION\\
Gold answer: GOLD\_ANSWER\\
Generated answer: GENERATED\_ANSWER\\
First, provide a short (one sentence) explanation of your reasoning, then finish with CORRECT or WRONG. Do NOT include both CORRECT and WRONG in your response, or it will break the evaluation script.}
\end{mdframed}

\subsubsection{Self-instruct DMR dataset generation}
The DMR question/answer pairs were generated using the following prompt and the original MSC dataset:
Your task is to write a "memory challenge" question for a simulated dialogue between two users.
\begin{mdframed}
\small
\texttt{You get as input:\\
- personas for each user (gives you their basic facts)\\
- a record of an old chat the two users had with each other\\\\
Your task is to write a question from user A to user B that test's user B's memory.\\
The question should be crafted in a way that user B must have actually participated in the prior conversation to answer properly, not just have read the persona summary.\\
Do NOT under any circumstances create a question that can be answered using the persona information (that's considered cheating).\\
Instead, write a question that can only be answered by looking at the old chat log (and is not contained in the persona information).\\\\
For example, given the following chat log and persona summaries:\\\\
old chat between user A and user B\\
A: Are you into surfing? I'm super into surfing myself\\
B: Actually I'm looking to learn. Maybe you could give me a basic lesson some time!\\
A: Yeah for sure! We could go to Pacifica, the waves there are pretty light and easy\\
B: That sounds awesome\\
A: There's even a cool Taco Bell right by the beach, could grab a bite after
B: What about this Sunday around noon?\\
A: Yeah let's do it!\\\\
user A persona:\\
I like surfing\\
I grew up in Santa Cruz\\\\
user B persona:\\
I work in tech\\
I live in downtown San Francisco\\\\
Here's an example of a good question that sounds natural, and an answer that cannot be directly inferred from user A's persona:\\\\
User B's question for user A\\
B: Remember that one time we went surfing? What was that one place we went to for lunch called?\\
A: Taco Bell!\\\\
This is an example of a bad question, where the question comes across as unnatural, and the answer can be inferred directly from user A's persona:\\\\
User B's question for user A\\
B: Do you like surfing?\\
A: Yes, I like surfing\\\\
Never, ever, ever create questions that can be answered from the persona information.
}
\end{mdframed}

\subsubsection{Document Analysis Instructions}
Example instructions used in the preprompt for document analysis tasks.
\begin{mdframed}
\small
\texttt{You are MemGPT DOC-QA bot. Your job is to answer questions about documents that are stored in your archival memory. The answer to the users question will ALWAYS be in your archival memory, so remember to keep searching if you can't find the answer. Answer the questions as if though the year is 2018. }
\end{mdframed}

Questions were provided to \ours with the following prompt:  
\begin{mdframed}
\small
\texttt{Search your archival memory to answer the provided question. Provide both the answer and the archival memory result from which you determined your answer. Format your response with the format 'ANSWER: [YOUR ANSWER], DOCUMENT: [ARCHIVAL MEMORY TEXT]. Your task is to answer the question:}
\end{mdframed}
For baselines, the following prompt along with a retrieved list of documents was provided: 
\begin{mdframed}
\small
\texttt{Answer the question provided according to the list of documents below (some of which might be irrelevant. In your response, provide both the answer and the document text from which you determined the answer. Format your response with the format 'ANSWER: <YOUR ANSWER>, DOCUMENT: [DOCUMENT TEXT]'. If none of the documents provided have the answer to the question, reply with 'INSUFFICIENT INFORMATION'. Do NOT provide an answer if you cannot find it in the provided documents. Your response will only be considered correct if you provide both the answer and relevant document text, or say 'INSUFFICIENT INFORMATION'. Answer the question as if though the current year is 2018.}
\end{mdframed}

\subsubsection{LLM Judge (document analysis)}
In order to both check the correctness of the answer for the document analysis task, and also to ensure that the answer was properly derived from the provided text (rather than from the model weights), we used an LLM judge. The LLM judge was provided the answers generated by both baseline approaches and \ours, and asked to make a judgement with the following prompt: 
\begin{mdframed}
\small
\texttt{Your task is to evaluate whether an LLM correct answered a question. The LLM response should be the format "ANSWER: [answer], DOCUMENT: [document\_text]" or say "INSUFFICIENT INFORMATION". The true answer is provided in the format "TRUE ANSWER:[list of possible answers]". The questions is provided in the format "QUESTION: [question]". If the LLM response contains both the correct answer and corresponding document text, the response is correct. Even if the LLM's answer and the true answer are slightly different in wording, the response is still correct. For example, if the answer is more specific than the true answer or uses a different phrasing that is still correct, the response is correct. If the LLM response if "INSUFFICIENT INFORMATION", or the "DOCUMENT" field is missing, the response is incorrect. Respond with a single token: "CORRECT" or "INCORRECT".}
\end{mdframed}

\subsubsection{K/V Task Instructions}
The \ours agent was defined with the following persona, designed to encourage \ours to iteratively search: 
\begin{mdframed}
\small
\texttt{You are MemGPT DOC-QA bot. Your job is to answer questions about documents that are stored in your archival memory. The answer to the users question will ALWAYS be in your archival memory, so remember to keep searching if you can't find the answer. DO NOT STOP SEARCHING UNTIL YOU VERIFY THAT THE VALUE IS NOT A KEY. Do not stop making nested lookups until this condition is met.}
\end{mdframed}
Baselines were instructed with the following prompt: 
\begin{mdframed}
\small
\texttt{Below is a JSON object containing key-value pairings, all keys and values are 128-bit UUIDs, and your task is to return the value associated with the specified key. If a value itself is also a key, return the value of that key (do a nested lookup). For example, if the value of 'x' is 'y', but 'y' is also a key, return the value of key 'y'.}
\end{mdframed}







